\documentclass{article}
\usepackage[T1]{fontenc}
\usepackage{lmodern}
\usepackage{graphicx}
\usepackage{amsmath,amssymb}
\usepackage[a4paper,margin=2cm]{geometry}
\usepackage[absolute,overlay]{textpos}
\setlength{\TPHorizModule}{1mm}
\setlength{\TPVertModule}{1mm}
\usepackage[indonesian]{babel}
\usepackage{hyperref}
\setlength{\emergencystretch}{2em}
% unit: Halaman 86

\begin{document}

% === Halaman 86 (llm) ===

\section*{2 fmt Chunk (Format Chunk)}

Menjelaskan bagaimana audio disimpan.

Struktur:

\begin{tabular}{lll}
\textbf{Field} & \textbf{Ukuran} & \textbf{Keterangan} \\ \hline
"fmt " & 4 byte & ID chunk \\
Subchunk1Size & 4 byte & Biasanya 16 untuk PCM \\
AudioFormat & 2 byte & 1 = PCM \\
NumChannels & 2 byte & 1=Mono, 2=Stereo \\
SampleRate & 4 byte & 44100, 48000, dll \\
ByteRate & 4 byte & $\text{SampleRate} \times \text{NumChannels} \times \text{BitsPerSample}/8$ \\
BlockAlign & 2 byte & $\text{NumChannels} \times \text{BitsPerSample}/8$ \\
BitsPerSample & 2 byte & 16, 24, 32 \\
\end{tabular}

\vspace{20mm}

\begin{flushright}
\textbf{Contoh umum:}
\begin{itemize}
    \item PCM
    \item 2 channel (stereo)
    \item 44100 Hz
    \item 16 bit
\end{itemize}

Sumber: ChatGPT
\end{flushright}

\end{document}
